Alle für die Studierenden relevanten Veranstaltungen der Hochschule sind frei zugänglich.

Alle Lehrenden bieten während ihrer Lehrveranstaltungen und der Betreuung von Prüfungsleistungen eine wöchentliche Sprechstunde an oder zumindest die Möglichkeit einen Termin innerhalb einer Woche zu vereinbaren. Sofern die Lehrperson für mehr als eine Woche nicht erreichbar ist, muss eine Alternative kommuniziert werden.

Die Lehrenden bzw. Betreuenden sind fachlich und didaktisch kompetent und sind um eine gute Lehre bemüht.

Es werden nicht nur Lösungen reproduziert, sondern auch Lösungsstrategien vermittelt.

Erwartete Vorkenntnisse und Ziele der Veranstaltung werden den Studierenden im Vorhinein klar kommuniziert.

Im geplanten Zeitaufwand für eine Veranstaltung wird eine angemessene Selbstlernzeit berücksichtigt.

Alle Veranstaltungen werden jedes Semester evaluiert, sofern die Anonymität der Befragten gewährleistet bleibt. Die Ergebnisse sind für die Studierenden in ausgewerteter Form zugänglich und werden von der Lehrperson genutzt, um die Lehre zu verbessern.

Nach dem bestandenen ersten Studienabschnitt wird davon ausgegangen, dass alle Studierenden sich auf etwa gleichem Niveau befinden. Hierbei wird auch auf Schwankungen bei den Vorkenntnissen der Studierenden eingegangen, d.h. das erreichte Niveau ist unabhängig vom Zeitpunkt des Studienbeginns. Eventuell vorhandene und erkannte Mängel der Studierenden werden durch zusätzliche Veranstaltungen oder Hilfestellungen, wie z. B. Übungen, Zusatzmaterial ausgeglichen.

Für die Mehrheit der Studierenden genügen 40 Stunden Arbeitsaufwand pro Woche, um das Studium in Regelstudienzeit erfolgreich abzuschließen (bei mind. sechs freien Wochen im Jahr). Der Arbeitsaufwand beinhaltet sowohl die Zeit für den Besuch von Veranstaltungen als auch für die Vor- und Nachbereitung, Hausaufgaben, schriftliche Arbeiten und ähnliches.

\section{Vorlesung}

Eine Vorlesung ist eine regelmäßige und fortlaufende Unterrichtsveranstaltung, die von Professor*innen, Lehrbeauftragten oder wissenschaftlichen Mitarbeitenden im Vortragsstil gehalten wird. Ziel von Vorlesungen ist die Vermittlung fachlichen Wissens auf theoretischer Basis.

Der Lehrstoff ist inhaltlich und visuell so aufbereitet, dass die Studierenden in der Regel mehrheitlich nicht überfordert sind. Dazu wird der Vortrag fachlich korrekt und sprachlich gut verständlich gehalten und ist didaktisch hochwertig. Eine sich durch das gesamte Semester ziehende Struktur des Lehrstoffes ist klar von den Studierenden erkennbar.

Zur Klärung fachlicher Fragen während der Veranstaltung ist ein gewisses Maß an Interaktivität gegeben. Hierbei werden Thematik und Gruppengröße berücksichtigt.

Durch Bereitstellung und/oder Verweise auf begleitende Lehrmaterialien ist es den Studierenden möglich, das Lernziel auch autodidaktisch zu erreichen sowie in der Vorlesung angeeignetes Wissen weiter zu vertiefen. Um einen hohen Vernetzungsgrad zwischen den Vorlesungen zu erreichen, gibt es fachliche Einordnungen der Themen und Ausblicke auf weiterführende Veranstaltungen.

Eine Vorlesung wird bei Grundlagenveranstaltungen von Übungen und/oder Tutorien begleitet. Das Verhältnis der Stundenzahl von Übungen/Tutorien zur Vorlesung beträgt mindestens 1:2.

\section{Übung / Tutorium}

Eine Übung bzw. ein Tutorium ist eine Kleingruppe, die von geeigneten Lehrverantwortlichen betreut wird und notwendigen Stoff und Übungsaufgaben behandelt. In einer Übung bzw. einem Tutorium wird die in der Vorlesung vermittelte Theorie angewandt und wiederholt, sowie erlernter Stoff gefestigt. Übungen und Tutorien beschäftigen sich mit der Konstruktion von Beispielen und Lösungen von Aufgabenstellungen und sind mit den zugehörigen Vorlesungen eng verknüpft.

Die Übungsaufgaben zu den Grundlagenvorlesungen werden korrigiert und kooperativ gelöst, während es bei anderen Vorlesungen akzeptabel ist auf vorhandene Lösungen zu verweisen und die autodidaktischen Fähigkeiten der Studierenden zu fordern und fördern.

Im Großteil der Zeit sollte die Mehrheit der Studierenden in der Lage sein der Übung zu folgen und aktiv mitzuarbeiten; der Schwerpunkt liegt auf Interaktivität.

\section{Seminar}

In einem Seminar tragen Studenten über ein vorher eigenständig aufbereitetes Thema vor. Dieses wird von fachlich kompetentem Lehrpersonal betreut. Ziel eines Seminars ist es, das eigenständige wissenschaftliche Arbeiten zu fördern und zur Präsentation von Ergebnissen zu befähigen. Die Studierenden entwickeln hierbei ein tiefergehendes fachliches Verständnis.

Alle Vorträge beziehen sich auf ein vorher bekanntgegebenes Rahmenthema.

Falls vorhanden ist unterschiedlicher Arbeitsaufwand vor Vergabe der Vorträge bekannt und auf Wunsch anzugleichen. Der Anspruch der Vortragsthemen korreliert mit der zur Verfügung stehenden Bearbeitungszeit, die in der Regel mindestens zwei Wochen beträgt. Während der Erarbeitungsphase steht eine kompetente Ansprechperson für Rückfragen zur Verfügung.

Die Teilnehmenden erhalten nach ihrem Vortrag Feedback vom Lehrpersonal sowie auf Wunsch auch vom Auditorium.
